\section{Conclusion}
\label{sec:conclusion}

IVAD frames agent-failure diagnosis as a selective decision over
machine-checkable causal claims rather than an unconstrained explanation. It
combines re-addressable trajectory evidence, deterministic eligibility,
controlled semantic verification, bounded revision and acquisition, and a
group-selective acceptance protocol. Under a frozen pipeline, independently
sampled preregistered groups, and simultaneous finite-candidate bounds, the
protocol provides a calibration-only PAC statement with an explicit
no-feasible-point outcome.

SpanVouch realizes this trust boundary as a versioned, recoverable open-source
system. On the frozen 36-candidate benchmark, deterministic verification
accepted all 20 valid reports and intercepted all 16 injected defects. The
local v0.7.0 Windows candidate run recorded 1,803 passes, one skip, three known
CRLF-only frozen-reference failures, and 92.14\% coverage; it is therefore not
a fully green release gate. The 24-cell offline matrix reproduced the complete
multi-domain, multi-framework artifact path with zero provider and GPU calls.
The passing checks establish executable contract behavior and exercise the
system and reproducibility paths. The v0.7 formal bundle additionally records
DeepSeek-only observations: 2,148 scheduled/evaluated plans, zero missing cells,
completed B0--B3, and policy-skipped B4/B5 with no Qwen results. The observed B3-versus-B2
risk and coverage differences describe a risk--coverage tradeoff in this run;
they are not causal, cross-model, or certificate claims. H1--H5 remain
unresolved. Together, IVAD and SpanVouch provide a precise protocol and an
engineering substrate for auditable acceptance of agent-failure diagnoses while
keeping unsupported claims fail-closed.
